\begin{revadded}
\subsection{Output objective}

For an input component $S$, the goal is to retain the unknown repair-relevant endpoints in a compact subgraph $H\subseteq S$. We measure retention and compression separately. The implementation uses a capped evidence selector; it does not solve a coverage--compression optimization problem with access to the hidden endpoint labels.

\subsection{SCC-aware local rollouts}

Tarjan decomposition identifies each non-trivial \scc{} \citep{tarjan1972depth}. Within it, \samcgs{} acquires evidence using windows of at most four records. A window starts from a seed and grows through incoming or outgoing frontier edges; a scheduled relation probe can instead revisit an implicated pair. If the frontier is exhausted, the implementation may add one remaining record. Thus windows are bounded expansions, not complete radius neighborhoods.

Node and edge statistics are keyed by physical records and directed dependencies. For a visited seed or frontier edge $x$, the base score is
\begin{equation}
 \begin{aligned}
 A_t(x)={}&\bar r_t(x)+c_t\sqrt{\frac{\log(1+N_t)}{n_t(x)+\ell_t(x)+\epsilon}}\\
       &+\frac{\lambda\pi_t(x)}{1+\delta\sqrt{n_t(x)}}.
 \end{aligned}
 \label{eq:ucb}
\end{equation}
Here $\bar r_t$ is mean node risk or edge conflict frequency, $n_t$ counts observations, and $\ell_t$ is weighted virtual loss for concurrent workers. $N_t$ sums node visits for seed selection or candidate-edge visits for expansion. Selection mixes $A_t$ with normalized Gamma noise: $(1-d)A_t+d\,g\xi_t$, with $g=10$ for seeds and $5$ for edges; seeds with zero visits and zero weighted virtual loss receive infinite priority. This combines UCB exploration with visit-decayed evidence bias \citep{kocsis2006bandit,chaslot2008progressive}.

The node prior normalizes a non-negative weighted sum of repair-entry, local-subgraph, conflict-endpoint, AMAF-style, path-fragment, conflict-count, and discovery signals. Edge priors additionally use critical-pair, semantic, explicit-conflict, reverse-edge, and conflict-frequency evidence. Appendix~\ref{app:priors} lists their fixed coefficients; the frontier prior also adds $0.25$ times the entering node's prior, capped at one.
\begin{equation}
 \pi_t^{\mathrm{base}}(x)=\frac{w^\top\phi_t(x)}{\max_{x'}w^\top\phi_t(x')},
 \label{eq:relation-prior}
\end{equation}
with zero prior when the denominator is zero. These heuristics reuse evidence across rollouts, in the spirit of progressive bias and AMAF/RAVE \citep{gelly2007combining,browne2012survey}.

\subsection{Accumulated evidence and discovery signal}

The local evaluator reads visible record text and directed edges and returns node risk scores, explicit conflicts, local-subgraph nodes, repair entries, and optional critical-pair features. A transposition table reuses evaluations of the same sorted node set up to its configured reuse limit. Every consumed evaluation updates node running averages, pair-support counts, relation reasons, and a persistent critical-pair ledger. Pair hypotheses are keyed by unordered endpoint pairs; directed edges retain separate exploration statistics. There is no exponential-moving-average relation update.

The legacy \oc{} signal is a heuristic node alarm. After warm-up, a node is flagged when at least two conditions hold: its mean risk is more than 1.5 pooled standard deviations above the pooled mean, its average empirical percentile exceeds 0.8, and more than 80\% of its scores exceed the pooled median. Flags persist. This rule uses node-score histories, not a quantile of maximum relation conflicts, and supplies neither conformal coverage nor calibrated conflict probabilities. Appendix~\ref{sec:appendix_formula_details} specifies the updates.

\subsection{Final core selection}
\label{sec:dynamic-core}

After search, the selector reconstructs a core from all accumulated evidence, under cap $B(S)=\min(|S|,\max(4,\lceil0.45|S|\rceil))$. Let $e_T(v)$ be the mean over channels that are nonzero somewhere in the component, each normalized within that component and $b_T(v)$ the strongest qualifying pair score touching $v$. Candidate priority is
\begin{equation}
 s_T(v)=e_T(v)+0.24b_T(v)+0.16\mathbf{1}[v\in O_T],
 \label{eq:priority}
\end{equation}
where $O_T$ is the set of flagged nodes. Starting from an empty selection, the rule seeds protected critical pairs and strong relations, adds nodes in priority order, applies absolute/relative-score and gap stopping rules, then closes supported pairs within the cap. It is a heuristic, not an argmax solver with a greedy approximation guarantee.

``Dynamic'' refers to reselecting from accumulated evidence rather than freezing the first admitted nodes. The production loop constructs its final core once, after search; a core reconstructed at an earlier prefix does not feed back into future search. Algorithm~\ref{alg:samcgs} makes this ordering explicit. Section~\ref{sec:component_tests} compares final reselection with early locking and a simpler terminal selector on identical traces.

\begin{algorithm}[t]
\small
\caption{\samcgs{} for one strongly connected component}
\label{alg:samcgs}
\begin{algorithmic}[1]
\REQUIRE Component $S$, rollout budget $T$, fixed configuration
\STATE Initialize node/edge statistics, pair ledger, and cache
\FOR{$t=1$ to $T$}
 \STATE Select a relation probe or UCB seed/frontier window
 \STATE Reuse a cached evaluation or query the local evaluator
 \STATE Update node means, pair counts, reasons, and ledger
 \STATE Update the persistent node-discovery flags
\ENDFOR
\STATE Aggregate evidence and select the capped final core $H$
\RETURN $H$ and ranking by mean node risk $\bar q_T(v)$
\end{algorithmic}
\end{algorithm}

\subsection{Computational feasibility}

Tarjan preprocessing costs $O(|V|+|E|)$. For each component, at most $T$ local evaluations are needed before retries, with window text bounded by four records; cache hits reduce new calls. Evidence storage includes node/edge statistics, observed pair hypotheses, and cached windows. Observed pairs can include non-edges, so memory is not bounded by the original edge count alone. Prefix readouts support anytime inspection \citep{zilberstein1996anytime}; they are analyzed in Appendix~\ref{app:budget_prefix}. Section~\ref{sec:cost_analysis} reports measured inference costs rather than treating rollout count as token cost.
\end{revadded}
